% !PRINTED 318
\weberpage{318}

Um eine obere Grenze für den absoluten Werth der imaginären Wurzeln, der
auch für den Fall complexer Co\"efficienten noch gültig ist, zu finden,
betrachten wir die Function
\[
f(x)=x^n+a_1x^{n-1}+a_2x^{n-2}+\cdots+a_n,
\]
worin die Co\"efficienten $a_1, a_2\ldots a_n$ reell oder imaginär sind.
Wir bezeichnen den absoluten Werth einer complexen Grösse $a$ wie früher
durch $|a|$ und haben dann, wenn wir annehmen, der absolute Werth von
$x$ sei grösser als 1,
\[
\left| a_1+\frac{a_2}{x}+\cdots+\frac{a_n}{x^{n-1}} \right| <
|a_1|+|a_2|+\cdots+|a_n|=\alpha.
\]
Nehmen wir also $x$ so an, dass sein absoluter Werth grösser ist als
jede der beiden Zahlen 1 und $\alpha$, so kann der Ausdruck
\[
x^{1-n}f(x)=x+a_1+\frac{a_2}{x}+\cdots+\frac{a_n}{x^{n-1}}
\]
nicht verschwinden, da der absolute Werth von $x$ grösser ist als der
absolute Werth der Summe aller übrigen Glieder der rechten Seite.

Wenn wir also, je nachdem $\alpha$ grösser oder kleiner als 1 ist, einen
Kreis mit dem Radius $\alpha$ oder 1 um den Nullpunkt als Mittelpunkt
legen, so liegen ausserhalb dieses Kreises gewiss keine Wurzeln von
$f(x)$ mehr.

\begin{center}
\S.\ 105.

\gesperrt{Abschätzung der imaginären Wurzeln.}
\end{center}

Wir haben im \S.\,96 gesehen, wie wir aus der Charakteristik eines
gewissen Curvensystems die genaue Anzahl der complexen Wurzeln einer
Gleichung bestimmen können, die im Inneren einer geschlossenen Curve
liegen. Auch dies Kennzeichen lässt sich vereinfachen, wenn man nicht
mehr die genaue Anzahl der Wurzeln, sondern nur eine obere Grenze für
diese Zahl finden will, was in vielen Fällen auch zur genauen
Bestimmung ausreicht.

Es sei also wie in \S.\,96
\weq{1}{F(x)=\varphi(x,y)+i\psi(x,y)}
eine reelle oder complexe Function des complexen Arguments
%CONT>>
